% Per-competition case studies -- companion document.
% Compile with pdfLaTeX (Overleaf default). No external assets required.
\documentclass[11pt,a4paper]{article}

\usepackage[margin=2.5cm]{geometry}
\usepackage{amsmath,amssymb}
\usepackage{newtxtext,newtxmath}
\usepackage{float}
\usepackage{booktabs}
\usepackage{caption}
\usepackage{microtype}
\usepackage{array}
\usepackage{enumitem}
\usepackage[hyphens]{url}
\usepackage[colorlinks=true, linkcolor=black, citecolor=black, urlcolor=blue!50!black]{hyperref}
\usepackage{xcolor}

\emergencystretch=1.5em
\captionsetup{font=small,labelfont=bf}

\title{\textbf{Per-Competition Case Studies}\\
\large Companion to \emph{Benchmarking AI Agents for Kaggle Competitions}\\
\normalsize\textcolor{orange!85!black}{All printed numbers are real measurements as of 2026-07-31.}}
\author{Wei-Hao Huang\\\small Advisor: Tso-Jung Yen}
\date{Draft compiled July 31, 2026}

\begin{document}
\maketitle

\begin{center}\textbf{\large Why this document exists}\end{center}
\noindent The main report states its claims at the level of methods and evaluation sets; the
evidence behind several of them is a handful of competitions examined very closely. That
forensic depth belongs neither in the report's running text, where it pulls the center of
gravity toward single competitions, nor in the engineering log, which records what broke while
building the machinery. It lives here. Most of these cases concern s3e19, the benchmark's
canonical failure; the fourth records the s5e1 setting behind the pre-registered prediction.

\vspace{0.5em}

\section{The CV--Leaderboard Reversal Mechanism}

Detail behind the main report's analysis \emph{Local CV and the Leaderboard Disagree}. CV and
the leaderboard measure different things: CV asks ``how well do we do on resamples of the
training data,'' the leaderboard asks ``how well do we do on a batch of never-seen data.'' The
two correlate strongly when distributions are stable, and decouple as soon as the test--train
relationship changes. s3e19 is the extreme case --- the matched-split CVs sit at 10--14
(AIDE's own selected CV, 4.4, is the leaked shuffled-KFold number of the next case) while
their leaderboard scores sit at 48--50: everyone misfires together, by amounts unrelated to the
CV ordering, into the lower mode of a bimodal leaderboard that CV had no way to predict. The
bimodality is sharp in the percentiles: on the 1{,}174-team leaderboard the 10th percentile is
6.1 while the 25th percentile jumps to 22.8 --- almost no middle ground exists between the mode
that solved the one-year extrapolation (265 teams below 10) and the mode that did not (570
teams in 40--60).

\section{AIDE's Silent Split Switch, Step by Step}

Detail behind the main report's failure-mode analysis. AIDE's metric trajectory on s3e19 jumps
from 18.42 to 4.41 within a single step:

\begin{verbatim}
step 0-10 : 18.42 ... 22.20 ... 18.42   (time-based validation)
step 11   : 4.41                        <- switched to KFold(shuffle=True)
step 12-19: 4.40, 4.37, 4.40, 4.49      (nine steps polishing a leaked number)
\end{verbatim}

Step 11's plan was merely to fix a LightGBM API compatibility issue; while rewriting the
training loop it incidentally replaced the validation split with shuffled KFold. The SMAPE
definition is identical before and after. This study confirmed the effect with a controlled
experiment (same features, same model, split changed only): \texttt{KFold(shuffle=True)} yields
SMAPE 4.56, while training on 2017--2020 and validating on 2021 yields 20.41 --- a $4.5\times$
gap. The competition's test set is one full future year; shuffling rows lets the model
interpolate within sequences it has memorized. AIDE noticed nothing.

The contrast with my-agent is the sharpest method-level distinction the benchmark surfaced:
my-agent also tried shuffled KFold in the same competition (SMAPE 4.2814) but explicitly tagged
it in \texttt{experiments.json} as ``DIAGNOSTIC: not used for submission,'' its purpose being
to isolate whether a score regression came from the split or from the features.

\section{The Two Experience-Library Entries Behind the Midterm Correction}

Tracing why the midterm had over-attributed s3e19's lower mode to the missing GDP join, we
found the experience library held both readings at once: one entry correctly identified the
one-year total extrapolation as the dominant obstacle, while a neighbouring entry over-claimed
that joining GDP \emph{dissolves} it. That both coexisted --- and that we acted on the
optimistic one --- is itself a finding: the library's weak link is retrieval, not content. The
over-claiming entry has since been corrected against the leaderboard evidence reported in the
main report's injection-layer experiments.

\section{s5e1 (``Forecasting Sticker Sales''): the Setting Behind the Pre-Registered Prediction}

playground-series-s5e1 is a fourth firing-class competition in the same
country$\times$store$\times$product panel family as s3e19 and tps-sep-2022, with six real
countries and a \emph{3-year} test horizon against the predecessors' one year. The
proportionality diagnostic read it as the worst violation in the class: cross-country
dispersion $0.05$--$0.21$, with the proportionality relation broken in 6 of 7 training years.
On that reading --- and before any model was fitted or anything submitted --- the prediction of
the main report's Result 2b (\texttt{ratio\_target} loses to \texttt{join\_feature}, by a
larger relative margin than on either predecessor) was committed to the competition's dossier
and to version control with its falsifier stated alongside. The three searched arms that then
killed it (private MAPE $0.15751$ / $0.15626$ / $0.12417$) applied no remedy for the six
violated years: the worst diagnosed violation in the class coincided with the ratio form's
largest win, which is what demoted the diagnostic from a gate to a recorded note.

\section{s3e19: the Residual Decomposition in Full}

The level/shape split reported in the main report's Result 3 comes from global rescalings of
the best arm's predictions, labelled as diagnostics and never reported as agent results:
$\times 1.0$ scores private SMAPE 48.497, $\times 1.6$ scores 32.932, and $\times 2.2$ scores
46.776 --- hence roughly 15.5 points of yearly-level error, with roughly 33 remaining as shape
against the leaderboard top of 4.67. The evidence that the needed $1.6$ multiplier is
unlearnable in sample: the proportionality relation implies $1.05$, and the evaluator's own
global-scale grid search settles on $1.0$. The transfer test of the experience library's other
candidate recipe: multiplicative panel decomposition, which had beaten a GBDT $4.19$ to $8.43$
on tps-jan-2022, scores $12.35$ on s3e19's held-out 2021 against the GBDT's $11.59$, and
$12.34$ even when handed the true held-out yearly total --- the recipe fails to transfer even
though its proportionality precondition appears satisfied on the training years.

\end{document}
